\chapter{Opis struktury projektu}
\label{cha:struktura}

W tym rozdziale opisano, jak zbudowana jest aplikacja: jakie technologie zostały wykorzystane, jak
zorganizowane są dane w bazie, jak wygląda hierarchia klas oraz jakie są minimalne wymagania
potrzebne do uruchomienia programu.

\section{Architektura i wykorzystane technologie}
\label{sec:architektura}

Aplikacja \emph{Testownik} to program desktopowy napisany w języku C\#~\cite{msdocsCSharp} na
platformie .NET. Jego
budowa jest modularna i opiera się na podziale na warstwy, co zapewnia czytelność kodu oraz
oddzielenie logiki od interfejsu użytkownika i od sposobu przechowywania danych. Wyróżnić można
cztery główne warstwy:

\begin{itemize}
  \item \textbf{Model (\texttt{Models})} - zbiór klas encji odwzorowujących rzeczywiste obiekty
        dziedziny (użytkownik, grupa, test, pytanie, podejście, odpowiedź). Klasy te są mapowane na
        tabele bazy danych przez Entity Framework Core.

  \item \textbf{Dostęp do danych (\texttt{Data})} - konfiguracja kontekstu bazy danych
        (\texttt{AppDbContext}), reguły mapowania relacji i dziedziczenia, mechanizm migracji oraz
        klasa zasilająca bazę danymi początkowymi (\texttt{DbSeeder}).

  \item \textbf{Logika i usługi (\texttt{Services})} - warstwa pośrednia zawierająca repozytorium
        operacji bazodanowych (\texttt{DbRepository}), serwis oceniania (\texttt{GradingService}),
        eksport do CSV (\texttt{CsvExporter}), haszowanie haseł (\texttt{PasswordHash}) oraz wspólną
        obsługę wyjątków (\texttt{Safe}). W tej warstwie znajdują się także obiekty transferu danych
        (DTO), oddzielające encje bazodanowe od warstwy widoku.

  \item \textbf{Interfejs użytkownika (\texttt{Views})} - ekrany aplikacji zbudowane w technologii
        WPF, odpowiedzialne za prezentację danych i obsługę interakcji z użytkownikiem. Widoki
        komunikują się z bazą wyłącznie poprzez warstwę usług.
\end{itemize}

Do budowy interfejsu graficznego wykorzystano technologię \emph{Windows Presentation Foundation}
(WPF)~\cite{msdocsWPF}, która pozwala na tworzenie nowoczesnych, stylizowanych okien aplikacji
desktopowych. Warstwę dostępu do danych zrealizowano przy użyciu biblioteki \emph{Entity Framework
Core}~\cite{msdocsEFCore} - nowoczesnego mechanizmu mapowania obiektowo-relacyjnego (ORM), który
umożliwia operowanie na bazie danych za pomocą obiektów języka C\#, bez konieczności ręcznego
pisania większości zapytań SQL. Jako silnik bazy danych zastosowano
\emph{PostgreSQL}~\cite{postgresql}, łączony z aplikacją za pośrednictwem dostawcy
\emph{Npgsql}~\cite{npgsql}. Do kontroli wersji kodu wykorzystano system Git.

\section{Zarządzanie danymi - baza danych}
\label{sec:bazaDanych}

Wszystkie informacje o użytkownikach, grupach, testach, pytaniach, podejściach i odpowiedziach są
przechowywane w relacyjnej bazie danych PostgreSQL. Schemat bazy jest tworzony i aktualizowany
automatycznie za pomocą migracji Entity Framework Core, które wersjonują kolejne zmiany struktury.
Najważniejsze tabele bazy danych przedstawiono w tabeli~\ref{tab:tabele}.

\begin{table}[H]
  \centering
  \caption{Najważniejsze tabele bazy danych aplikacji.}
  \label{tab:tabele}
  \begin{tabularx}{\linewidth}{>{\bfseries}l X}
    \toprule
    Tabela & Przeznaczenie \\
    \midrule
    Users & Konta użytkowników (wspólna tabela nauczycieli i uczniów, rozróżnianych
            dyskryminatorem); login, skrót hasła, imię i nazwisko. \\
    Groups & Grupy uczniów (nazwa, opis). \\
    StudentGroups & Tabela łącząca uczniów z grupami (relacja wiele-do-wielu). \\
    Quizzes & Testy (nazwa, przedmiot, limit czasu, próg zaliczenia, liczba podejść). \\
    Questions & Pytania (wspólna tabela pytań zamkniętych i otwartych); treść, punkty, kolejność. \\
    AnswerOptions & Warianty odpowiedzi pytań zamkniętych; treść, flaga poprawności. \\
    TestGroups & Tabela łącząca testy z grupami (relacja wiele-do-wielu). \\
    QuizAttempts & Podejścia uczniów do testów; wynik, procent, status, czas, data. \\
    StudentAnswers & Odpowiedzi ucznia (wspólna tabela odpowiedzi zamkniętych i otwartych). \\
    SelectedOptions & Warianty zaznaczone przez ucznia w odpowiedzi zamkniętej. \\
    \bottomrule
  \end{tabularx}
\end{table}

Między tabelami zachodzą relacje jeden-do-wielu (np. test posiada wiele pytań, podejście posiada
wiele odpowiedzi) oraz dwie relacje wiele-do-wielu, zrealizowane przez tabele łączące: uczeń może
należeć do wielu grup, a grupa zrzeszać wielu uczniów (\texttt{StudentGroups}); analogicznie test
może być przypisany do wielu grup, a grupa mieć wiele testów (\texttt{TestGroups}). Integralność
danych zapewniają klucze obce wraz z regułami usuwania: kaskadowymi tam, gdzie dane podrzędne nie
mają sensu bez nadrzędnych (usunięcie testu usuwa jego pytania i podejścia), oraz ograniczającymi
tam, gdzie usunięcie mogłoby naruszyć już oddane odpowiedzi. Pełny schemat powiązań przedstawia
diagram ERD na rysunku~\ref{rys:erd}.

\figscreen{erd.png}{0.95}{Diagram ERD bazy danych aplikacji.}{rys:erd}

\section{Hierarchia klas i kluczowe metody}
\label{sec:hierarchiaKlas}

Projekt w szerokim zakresie wykorzystuje mechanizmy programowania obiektowego, w szczególności
dziedziczenie. W modelu danych występują trzy hierarchie klas, mapowane na bazę danych strategią
\emph{Table-per-Hierarchy} (TPH), w której cała hierarchia jest zapisywana w jednej tabeli,
a konkretny typ rozróżnia kolumna dyskryminatora.

\begin{itemize}
  \item \textbf{Hierarchia użytkowników} - abstrakcyjna klasa \texttt{User} przechowuje wspólne
        dane (identyfikator, login, skrót hasła, imię i nazwisko) oraz definiuje metodę abstrakcyjną
        \texttt{GetRoleLabel()}. Dziedziczą po niej klasy \texttt{Teacher} oraz \texttt{Student}.
        Uczeń posiada dodatkowo kolekcję grup, do których należy, oraz kolekcję swoich podejść do
        testów.

  \item \textbf{Hierarchia pytań} - abstrakcyjna klasa \texttt{Question} przechowuje treść, liczbę
        punktów i kolejność pytania. Dziedziczą po niej \texttt{ClosedQuestion} (pytanie zamknięte,
        posiadające warianty odpowiedzi oraz znacznik wielokrotnego wyboru) i \texttt{OpenQuestion}
        (pytanie otwarte z odpowiedzią wzorcową). Właściwość \texttt{RequiresManualGrading}
        rozróżnia, czy pytanie wymaga ręcznej oceny.

  \item \textbf{Hierarchia odpowiedzi} - abstrakcyjna klasa \texttt{StudentAnswer} jest bazą dla
        \texttt{ClosedAnswer} (przechowuje zaznaczone warianty) oraz \texttt{OpenAnswer}
        (przechowuje wpisany tekst). Pozwala to trzymać w jednej kolekcji odpowiedzi obu typów.
\end{itemize}

Oprócz dziedziczenia w projekcie wykorzystano interfejs \texttt{IPublishable}, implementowany przez
klasę \texttt{Quiz}. Definiuje on właściwość \texttt{IsPublished} oraz metody \texttt{Publish()}
i \texttt{Unpublish()}, służące do sterowania widocznością testu dla uczniów (stan aktywny
i nieaktywny).

Logikę aplikacji realizują klasy warstwy usług:

\begin{itemize}
  \item \texttt{DbRepository} - centralne repozytorium operacji na bazie danych. Hermetyzuje
        tworzenie kontekstu, pobieranie i zapisywanie testów, grup, podejść i wyników oraz mapowanie
        encji na obiekty transferu danych. Udostępnia m.in. metody zapisu nowego podejścia
        (\texttt{AddAttempt}), ręcznej aktualizacji oceny (\texttt{UpdateGrading}), pobrania
        najlepszego wyniku ucznia (\texttt{GetBestResult}), wyboru testów danego nauczyciela
        (\texttt{GetTestsForTeacher}), zmiany widoczności testu (\texttt{SetTestActive}) oraz
        bezpiecznego usuwania testu (\texttt{DeleteTest}).

  \item \texttt{GradingService} - statyczna klasa odpowiedzialna za ocenianie. Metoda
        \texttt{Grade} wyznacza liczbę punktów za pojedyncze pytanie (z punktami częściowymi
        w pytaniach wielokrotnego wyboru), a metoda \texttt{Recompute} przelicza cały wynik
        podejścia, uwzględniając pytania oczekujące na ręczną ocenę.

  \item \texttt{CsvExporter} - eksportuje wyniki testu do pliku CSV z użyciem strumieni
        (\texttt{StreamWriter}), dbając o poprawne kodowanie polskich znaków oraz cytowanie wartości.

  \item \texttt{PasswordHash} - wylicza skrót SHA-256 hasła, dzięki czemu hasła nie są
        przechowywane w bazie w postaci jawnej.

  \item \texttt{Safe} - wspólny mechanizm obsługi wyjątków; wykonuje operacje bazodanowe
        i wejścia/wyjścia w bloku zabezpieczającym, prezentując użytkownikowi czytelny komunikat
        zamiast powodować awarię aplikacji.
\end{itemize}

Wzajemne powiązania najważniejszych klas modelu przedstawia diagram UML na
rysunku~\ref{rys:uml}.

\figscreen{uml.png}{0.95}{Diagram UML klas modelu aplikacji.}{rys:uml}

\section{Minimalne wymagania programu}
\label{sec:wymaganiaMin}

\begin{itemize}
  \item \textbf{System operacyjny:} Windows 10 lub Windows 11.
  \item \textbf{Środowisko uruchomieniowe:} platforma .NET (wersja zgodna z projektem) wraz
        ze wsparciem WPF.
  \item \textbf{Serwer bazy danych:} PostgreSQL (zalecana wersja 13 lub nowsza), działający lokalnie
        lub zdalnie, z dostępem do bazy aplikacji.
  \item \textbf{Dodatkowe narzędzia:} środowisko programistyczne wspierające .NET (np. Visual
        Studio), system kontroli wersji Git oraz narzędzie do zarządzania bazą danych (np. pgAdmin).
\end{itemize}

\section{Uruchomienie aplikacji}
\label{sec:uruchomienie}

Aplikacja łączy się z bazą danych PostgreSQL przy użyciu parametrów zdefiniowanych w klasie
\texttt{AppDbContext}: serwer \texttt{localhost}, port \texttt{5432}, baza \texttt{testownik},
użytkownik \texttt{postgres}. Przed pierwszym uruchomieniem należy zapewnić działający serwer
PostgreSQL z dostępem do tych parametrów.

Schemat bazy danych zakładany jest automatycznie. Przy starcie aplikacja wywołuje metodę
\texttt{Database.Migrate()}, która tworzy bazę i nakłada wszystkie migracje Entity Framework Core,
a następnie klasa \texttt{DbSeeder} wypełnia bazę danymi początkowymi (przykładowy nauczyciel,
uczniowie, grupy oraz testy). Ręczne tworzenie tabel nie jest więc wymagane.

Kolejne kroki uruchomienia:
\begin{enumerate}
  \item uruchomić serwer PostgreSQL i upewnić się, że jest dostępny na \texttt{localhost:5432};
  \item otworzyć projekt w środowisku Visual Studio i uruchomić aplikację (lub wykonać polecenie
        \texttt{dotnet run} w katalogu projektu);
  \item przy pierwszym uruchomieniu baza zostanie utworzona i zasilona danymi demonstracyjnymi;
  \item zalogować się przy użyciu jednego z kont demonstracyjnych.
\end{enumerate}

Konta demonstracyjne (hasło dla wszystkich: \texttt{demo1234}):
\begin{itemize}
  \item nauczyciel - login \texttt{m.dabrowski};
  \item uczniowie - loginy \texttt{anna.k}, \texttt{piotr.n}, \texttt{maria.w} oraz kolejni.
\end{itemize}
